\documentclass[11pt]{article}

\usepackage[margin=1in]{geometry}
\usepackage{amsmath,amssymb,mathtools,bm}
\usepackage{booktabs,longtable,array}
\usepackage[hidelinks]{hyperref}

\newcommand{\E}{\mathbb{E}}
\newcommand{\Cov}{\operatorname{Cov}}
\newcommand{\Var}{\operatorname{Var}}
\newcommand{\Law}{\operatorname{Law}}
\newcommand{\diag}{\operatorname{diag}}
\newcommand{\rank}{\operatorname{rank}}
\newcommand{\tr}{\operatorname{tr}}
\newcommand{\one}{\bm{1}}
\newcommand{\zero}{\bm{0}}
\newcommand{\ThetaV}{\boldsymbol{\Theta}}

\title{SCI-NOI-002 Independent Mathematical Core\\
Finite ensembles, covariance calibration, and consumer validity}
\author{Independent Phase 1 audit derivation}
\date{Frozen against application parent\\
\texttt{d5015fe716971bf8ea617e8a187311bf5af05185}}

\begin{document}
\maketitle

\begin{abstract}
This document derives the conditions under which a finite conditional
randomization ensemble can support an empirical variance, covariance,
calibrated weight, signal-to-noise, significance, source-finding threshold,
feedback, aperture uncertainty, or production-count claim.  It distinguishes
uncentered second moments, centered randomization covariance, repeated-
physical-noise covariance, marginal inverse variance, full precision, and
consumer-specific effective precision.  It derives exact finite-stack scatter
identities for arbitrary cross-realization dependence; identifies the special
assumptions behind divisors $R$, $R-1$, weighted corrections, and analytic
inverse-covariance corrections; and states the extra distributional and
selection requirements for significance and adaptive consumers.

The result is a requirements core, not an implementation finding.  It selects
no estimator, calibration, threshold, feedback rule, aperture, production
realization count, or default.  In particular, 64 is optional validation
capacity only.  It is not a default, minimum, adequacy threshold, cap, or
universal requirement.
\end{abstract}

\section{Authority, quarantine, and exclusive scope}

The sole package-specific pre-core authority is the frozen
\texttt{SCI-NOI-002-XAUD-002} handoff, SHA-256
\texttt{9eb6c778409d344ce73387f44ac4a5429a89d43b8e768c19bf3da1ed6967c1e5}.
It admits the SCI-NOI-001 R3 conditional-sign ensemble contract at commit
\texttt{5a027c94ef9fc9c4a6e6cadc84af1c8a550d3508}, SHA-256
\texttt{27263ab3bf29ac8f098463455e540f13e783241a688ef2bc5cb15b1f2a4319da}.
R3 supplies the conditional ensemble law, its mean and covariance propagation,
finite realized assignments, source-imprint caveat, and use-specific target
projections.  It deliberately supplies no SCI-NOI-002 estimator or consumer
policy.

This derivation also conditions on the repository-level scientific contract:
the accepted MAP \texttt{weight\_I} quantity is a nonprecision gridding and
normalization coefficient.  Its inverse-squared signal unit is dimensional
metadata, not evidence that it is inverse variance or precision.  Raw science
validity, downstream numerical support, estimator validity, and consumer
validity remain separate.  MAP remains conditioned on PTC/VAL and
\texttt{existing\_use\_only}; JINC remains conditioned on SCI-MAP-002.  No
FLT, FRUIT, or other downstream package is reopened here.

No implementation source, package test, implementation diff, NOI-001 final
source audit, NOI-001 repair material, or sealed
\texttt{SCI-NOI-002-XAUD-001} evidence was used.  Consequently this document
makes no statement about what the application computes.  It gives the
independent equations against which a separately authorized trace can later
be performed.

\section{Conditional ensemble and target taxonomy}

Let $Y^{(r)}\in\mathbb R^P$ be one completed observation, coadd, or filtered
realization from the R3 \texttt{source\_imprinted\_current} ensemble, after a
fixed projection onto a declared common identity, unit, support, response, and
validity domain.  All equations are conditional on the exact realized state
$\ThetaV$ required by R3.  The realization index is $r\in\{1,\ldots,R\}$,
where $R$ is the completed count, not the requested or effective count.

For a one-realization marginal law define
\begin{align}
  \mu_r &= \E[Y^{(r)}\mid\ThetaV],
  \label{eq:marginal-mean}\\
  K_{rr} &= \Cov(Y^{(r)},Y^{(r)}\mid\ThetaV),
  \label{eq:marginal-covariance}\\
  M_{rr} &= \E[Y^{(r)}Y^{(r)\mathsf T}\mid\ThetaV]
          =K_{rr}+\mu_r\mu_r^{\mathsf T}.
  \label{eq:marginal-second-moment}
\end{align}
Across realized assignment positions define the joint-design covariance
\begin{equation}
  \Gamma_{rs}
  =\Cov(Y^{(r)},Y^{(s)}\mid\ThetaV),
  \qquad 1\le r,s\le R.
  \label{eq:cross-realization-covariance}
\end{equation}
Complement pairing, exact balance across realizations, orthogonal arrays,
sampling without replacement, duplicate assignments, reused streams, and
deterministic assignment lists can all make $\Gamma_{rs}\ne0$ for $r\ne s$.
The one-realization covariance and the cross-realization design are distinct
facts.

For a named use $u$, let $L_u$ be its fixed linear projection and
$\mathcal V_u$ its admitted validity domain.  The target need not be the
conditional randomization covariance.  Write
\begin{equation}
  T_u^\star
  =\Cov_{\star}(L_uY,L_uY\mid\ThetaV,\mathcal V_u)
  \label{eq:consumer-target}
\end{equation}
for the independently authorized target covariance.  The subscript $\star$
must say whether the target is conditional randomization, repeated physical
noise, an externally calibrated empirical law, or some other named law.
Unbiased estimation of the R3 conditional covariance does not by itself prove
$L_uK_{rr}L_u^{\mathsf T}=T_u^\star$.

At minimum, the following product meanings are distinct:
\begin{longtable}{>{\raggedright\arraybackslash}p{0.24\textwidth}
                  >{\raggedright\arraybackslash}p{0.69\textwidth}}
\toprule
Quantity & Required meaning \\
\midrule
\endhead
Uncentered ensemble second moment &
$M=\E[YY^{\mathsf T}]$ or its declared finite-stack analogue.  It equals
covariance only when the relevant mean is exactly zero. \\
Centered randomization covariance &
$K^\epsilon=\Cov_\epsilon(Y\mid\ThetaV)$ under the R3 assignment law, with
the fixed realized data and operators in the conditioning state. \\
Physical-noise covariance &
Covariance under repeated physical noise at a named pre-realization state.
R3 shows that this is not generally the conditional randomization covariance. \\
Marginal variance &
$v_p=K_{pp}$ for one declared pixel, coefficient, or projected statistic,
with squared signal unit. \\
Inverse marginal variance &
$w_p^{\rm marg}=1/v_p$ on a separately valid domain.  It is not a full
precision matrix and is biased when formed by naively inverting a noisy
variance estimate. \\
Precision matrix &
$\Pi=K^{-1}$, or a declared generalized/regularized inverse on a stated
subspace.  In general $\Pi_{pp}\ne1/K_{pp}$. \\
Consumer effective precision &
$w_u=[L_uKL_u^{\mathsf T}]^{-1}$ for a scalar consumer, or the corresponding
matrix on its projected domain.  It is tied to the exact operator $L_u$. \\
MAP coefficient &
The admitted nonprecision gridding/normalization \texttt{weight\_I}.  It is
none of the preceding empirical-noise quantities without separate authority. \\
\bottomrule
\end{longtable}

\section{Exact finite-stack identities}

\subsection{Equal-weight empirical centering}

Define
\begin{align}
  \bar Y &= \frac1R\sum_{r=1}^R Y^{(r)},
  \label{eq:sample-mean}\\
  S_R &= \sum_{r=1}^R
    (Y^{(r)}-\bar Y)(Y^{(r)}-\bar Y)^{\mathsf T}.
  \label{eq:centered-scatter}
\end{align}
Let $\bar\mu=R^{-1}\sum_r\mu_r$.  Direct expansion gives
\begin{align}
  \E[S_R\mid\ThetaV]
  ={}&\sum_{r=1}^R\Gamma_{rr}
     -\frac1R\sum_{r=1}^R\sum_{s=1}^R\Gamma_{rs}
  \nonumber\\
  &+\sum_{r=1}^R
    (\mu_r-\bar\mu)(\mu_r-\bar\mu)^{\mathsf T}.
  \label{eq:general-scatter-expectation}
\end{align}
Thus no divisor can be selected from $R$ alone.  Marginal means, marginal
covariances, and the complete cross-realization design determine the expected
scatter.

If all positions have common mean $\mu$ and common marginal covariance $K$,
Equation~\eqref{eq:general-scatter-expectation} reduces to
\begin{equation}
  \E[S_R\mid\ThetaV]
  =RK-\frac1R\sum_{r,s}\Gamma_{rs}.
  \label{eq:common-marginal-scatter}
\end{equation}
Only for independent assignments, $\Gamma_{rs}=0$ for $r\ne s$, does
\begin{equation}
  \E[S_R]=(R-1)K,
  \qquad
  \widehat K_{\rm iid}=\frac{S_R}{R-1}
  \label{eq:bessel-case}
\end{equation}
give the ordinary unbiased covariance estimator for $R\ge2$.  The divisor
$R$ instead gives the maximum-likelihood covariance under an independent
Gaussian model and has expectation $(R-1)K/R$ when the mean is estimated.
Calling either formula empirical variance without recording the mean and
design assumptions is incomplete.

For an exchangeable joint design with
\begin{equation}
  \Gamma_{rr}=K,
  \qquad
  \Gamma_{rs}=C\quad(r\ne s),
  \label{eq:exchangeable-design}
\end{equation}
the exact result is
\begin{equation}
  \E[S_R]=(R-1)(K-C).
  \label{eq:exchangeable-scatter}
\end{equation}
The $R-1$ estimator therefore targets $K-C$, not $K$.  If and only if a
separate authority establishes $C=\rho K$ with known scalar $\rho$, then
\begin{equation}
  \widehat K
  =\frac{S_R}{(R-1)(1-\rho)}
  \label{eq:scalar-correlation-correction}
\end{equation}
is an unbiased scalar correction when $\rho\ne1$.  A nonproportional $C$
cannot be repaired by replacing $R$ with one scalar effective count.  A known
joint design may instead define a matrix- or operator-valued correction on a
declared subspace, but its invertibility, conditioning, and bias are part of
the estimator contract.

Sampling from a finite assignment population is a concrete illustration.  Let
$\Omega$ contain $M_\Omega$ equally weighted maps, let $K_\Omega$ denote the
covariance of one uniform draw with population divisor $M_\Omega$, and draw
$R$ distinct assignments uniformly without replacement.  Then
\begin{equation}
  \Gamma_{rs}=-\frac{K_\Omega}{M_\Omega-1}
  \quad(r\ne s),
  \label{eq:without-replacement-cross-covariance}
\end{equation}
and the centered scatter obeys
\begin{equation}
  \E[S_R]
  =(R-1)\frac{M_\Omega}{M_\Omega-1}K_\Omega.
  \label{eq:finite-population-scatter}
\end{equation}
Thus $S_R/(R-1)$ is unbiased for the finite-population covariance defined with
divisor $M_\Omega-1$, while
$(M_\Omega-1)S_R/[M_\Omega(R-1)]$ targets the marginal uniform-draw covariance
$K_\Omega$.  Sampling with replacement, sampling without replacement, and
complete deterministic enumeration have different normalizations even when
their assignment set is the same.

The sampling covariance of the empirical mean is independently
\begin{equation}
  \Cov(\bar Y\mid\ThetaV)
  =\frac1{R^2}\sum_{r,s}\Gamma_{rs}.
  \label{eq:mean-covariance}
\end{equation}
For a fixed scalar projection $h_r=\ell^{\mathsf T}Y^{(r)}$ with common
marginal variance $v_h=\ell^{\mathsf T}K\ell$, a mean-specific information
descriptor is
\begin{equation}
  R_{\rm eff}^{\rm mean}(\ell)
  =\frac{R^2v_h}{\sum_{r,s}\ell^{\mathsf T}\Gamma_{rs}\ell},
  \label{eq:mean-effective-count}
\end{equation}
when the denominator is positive.  This can exceed $R$ for an antithetic
design or collapse to one for duplicates.  It does not give the effective
count for a variance, tail, covariance inverse, or source finder.

\subsection{Known mean, uncentered moments, and source imprint}

For a fixed reference mean $m_0$, define
\begin{equation}
  \widehat K(m_0)
  =\frac1R\sum_{r=1}^R
    (Y^{(r)}-m_0)(Y^{(r)}-m_0)^{\mathsf T}.
  \label{eq:known-mean-estimator}
\end{equation}
Its expectation is
\begin{equation}
  \E[\widehat K(m_0)]
  =\frac1R\sum_r\left[
    \Gamma_{rr}+(\mu_r-m_0)(\mu_r-m_0)^{\mathsf T}
  \right].
  \label{eq:known-mean-expectation}
\end{equation}
Cross-realization dependence affects its Monte Carlo uncertainty but not this
marginal expectation.  For a common marginal covariance and a truly known
$m_0=\mu$, it is unbiased for $K$ even when assignment positions are
dependent.  For $m_0=0$ it is the uncentered second moment.  R3 permits a
centered target sign law while a finite stack has a nonzero empirical sign
mean, and R3 also proves that deterministic source content can enter the
conditional second moment and covariance.  Consequently, choosing known zero
mean rather than empirical centering is a substantive target decision, not a
mere normalization preference.

\subsection{Unequal realization weights}

Let $a_r\ge0$, $\sum_ra_r=1$, and define
\begin{align}
  \bar Y_a &=\sum_ra_rY^{(r)},\\
  S_a &=\sum_ra_r(Y^{(r)}-\bar Y_a)
                    (Y^{(r)}-\bar Y_a)^{\mathsf T}.
  \label{eq:weighted-scatter}
\end{align}
With $\bar\mu_a=\sum_ra_r\mu_r$,
\begin{align}
  \E[S_a]={}&\sum_ra_r\Gamma_{rr}
  -\sum_{r,s}a_ra_s\Gamma_{rs}
  \nonumber\\
  &+\sum_ra_r(\mu_r-\bar\mu_a)(\mu_r-\bar\mu_a)^{\mathsf T}.
  \label{eq:weighted-scatter-expectation}
\end{align}
For independent identically distributed positions this becomes
\begin{equation}
  \E[S_a]=\left(1-\sum_ra_r^2\right)K,
  \qquad
  \widehat K_a=\frac{S_a}{1-\sum_ra_r^2}.
  \label{eq:weighted-iid-correction}
\end{equation}
The often quoted $R_w=(\sum_ra_r)^2/\sum_ra_r^2=1/\sum_ra_r^2$ is therefore
only a weight-concentration count in this case.  It does not correct unknown
cross-realization covariance or data-derived weights.  If $a_r$ is estimated
from the same realization stack, Equation~\eqref{eq:weighted-scatter-expectation}
does not apply conditionally without including the weight--map joint law.

\subsection{Missing realizations and changing support}

Every estimator must use completed realization identities.  If the valid set
depends on pixel or product, let $R_p$ be the count at pixel $p$ and $R_{pq}$
the common count for a covariance pair.  Substituting one global $R$ is then
wrong.  Pairwise deletion can produce a non-positive-semidefinite matrix
because different entries use different realization populations.  A valid
contract must choose one of: a fixed complete-case domain; a declared
missingness model and estimator; or separate products whose supports and
pairwise populations are explicit.  A missing required realization is not a
numerical zero, and a partially persisted stack cannot claim its requested
cardinality.

\section{Rank, inversion, and regularization}

For $P$ projected coordinates,
\begin{equation}
  \rank(S_R)\le\min(P,R-1),
  \label{eq:rank-bound}
\end{equation}
with a potentially smaller bound from duplicate, complement-equivalent, or
otherwise restricted even-moment assignments.  Thus a centered empirical
covariance cannot generically supply a full $P$-dimensional precision matrix
when $R\le P$.  A pseudoinverse then describes only the realized covariance
span; it is not full precision.  Binning, low-rank modeling, tapering,
shrinkage, or another regularizer changes the estimator and must name its
target, hyperparameter rule, bias, rank, and out-of-sample validation.

The independent Gaussian limit gives a useful falsification benchmark, not a
default model.  If $Y^{(r)}\stackrel{\rm iid}{\sim}\mathcal N_P(\mu,K)$ and
\begin{equation}
  \widehat K=\frac{S_R}{R-1},
  \label{eq:gaussian-sample-covariance}
\end{equation}
then $(R-1)\widehat K$ is Wishart with $R-1$ degrees of freedom.  For one
coordinate with variance $v$,
\begin{equation}
  \E[\widehat v^{-1}]
  =\frac{R-1}{R-3}\,v^{-1},
  \qquad R>3.
  \label{eq:inverse-variance-bias}
\end{equation}
Hence naive inversion is biased upward even when the variance estimator is
unbiased; $(R-3)\widehat v^{-1}/(R-1)$ is unbiased only in this exact
univariate iid Gaussian setting.  For a full $P$-dimensional inverse,
\begin{equation}
  \E[\widehat K^{-1}]
  =\frac{R-1}{R-P-2}\,K^{-1},
  \qquad R>P+2,
  \label{eq:inverse-covariance-bias}
\end{equation}
so $(R-P-2)\widehat K^{-1}/(R-1)$ is the corresponding analytic correction.
Replacing $R$ by a generic effective count in these formulae is unjustified:
correlated randomization designs need not be Wishart, and non-Gaussian sign
ensembles need not have Gaussian fourth moments.

Under the same iid Gaussian benchmark,
\begin{equation}
  \Cov(\widehat K_{ij},\widehat K_{k\ell})
  =\frac{K_{ik}K_{j\ell}+K_{i\ell}K_{jk}}{R-1},
  \label{eq:wishart-entry-covariance}
\end{equation}
and
\begin{equation}
  \frac{\operatorname{sd}(\widehat v)}{v}
  =\sqrt{\frac{2}{R-1}}.
  \label{eq:variance-relative-error}
\end{equation}
For a dependent ensemble, uncertainty of a variance estimate depends on the
cross-realization covariance of quadratic summaries
$q_r=(\ell^{\mathsf T}(Y^{(r)}-\mu))^2$, hence on fourth-order joint moments.
If the $q_r$ have common marginal variance $\tau_q^2$ and covariance matrix
$\Lambda_q$, the analogous descriptor
\begin{equation}
  R_{\rm eff}^{q}(\ell)
  =\frac{R^2\tau_q^2}{\one^{\mathsf T}\Lambda_q\one}
  \label{eq:quadratic-effective-count}
\end{equation}
applies only to the Monte Carlo mean of that exact quadratic summary when the
denominator is positive.  $R_{\rm eff}^{\rm mean}$ cannot substitute for it,
and neither effective count supplies a finite-$R$ bias correction by itself.

\section{Calibration and truthful weight identities}

Let $\widetilde K_u$ be a finite-stack estimate on the domain of consumer
$u$.  A calibration is an explicit, independently governed map
\begin{equation}
  \widehat T_u
  =\mathcal C_u(\widetilde K_u;D_{\rm cal},\Xi_u),
  \label{eq:calibration-operator}
\end{equation}
where $D_{\rm cal}$ identifies the calibration data and $\Xi_u$ records the
target law, response, region, estimator version, and fitting assumptions.  A
calibration claim requires evidence that $\widehat T_u$ estimates
$T_u^\star$, including overlap or independence between $D_{\rm cal}$, the
science numerator, and the ensemble used for $\widetilde K_u$.

A scalar scale calibration has the special form
\begin{equation}
  \widehat T_u=\widehat\alpha_u\widetilde K_u,
  \qquad \widehat\alpha_u>0.
  \label{eq:scalar-calibration}
\end{equation}
It can correct a scale functional for which proportionality has been
established.  It cannot repair a wrong correlation shape, source-imprint
morphology, support mismatch, response error, or tail law.  A spatially varying
calibration is a larger covariance model and needs correspondingly stronger
identifiability and validation.

Once a calibrated covariance is valid, three common weight constructions are
\begin{align}
  \widehat w_p^{\rm marg}
  &=\frac{1}{\widehat T_{u,pp}},
  \label{eq:marginal-inverse-variance}\\
  \widehat\Pi_u
  &=\widehat T_u^{-1},
  \label{eq:precision-estimate}\\
  \widehat w_u^{\rm eff}
  &=\frac{1}{\ell_u^{\mathsf}\widehat T_u\ell_u}
  \quad\text{for a scalar fixed projection }\ell_u.
  \label{eq:effective-consumer-weight}
\end{align}
Equation~\eqref{eq:marginal-inverse-variance} is an inverse marginal variance,
Equation~\eqref{eq:precision-estimate} is a full or declared regularized
precision, and Equation~\eqref{eq:effective-consumer-weight} is an effective
precision for one operator.  They are not interchangeable.  Each requires
finite positive denominator or a declared positive-definite subspace, and
each carries inverse squared output unit.  Invalid or unavailable estimates
use an explicit validity state; a zero, infinity, NaN, or negative number is
not an undocumented sentinel.

At least the following metadata is required before a quantity may be labeled
empirical variance, calibrated weight, or precision:
\begin{longtable}{>{\raggedright\arraybackslash}p{0.22\textwidth}
                  >{\raggedright\arraybackslash}p{0.71\textwidth}}
\toprule
Field class & Required fact \\
\midrule
\endhead
Target & Conditional randomization, repeated physical noise, calibrated null,
or another exact law; observation/coadd/filter stage and raw-parent identity. \\
Centering & Known reference mean, empirical mean, or uncentered second moment;
the mean estimate and whether deterministic source imprint is admitted. \\
Finite design & Requested/effective/completed $R$; joint assignment design;
duplicates, complements, unique even patterns, weighting, missingness, and the
exact normalization or design correction. \\
Covariance domain & Pixel, block, stationary kernel, low-rank subspace, full
matrix, or fixed consumer projection; units, shape, frame, indexing, response,
support, and validity. \\
Calibration & Calibration data identity; independence/overlap; scalar,
spatial, matrix, or distributional transform; estimator/version and its
uncertainty. \\
Regularization & None, pseudoinverse, shrinkage, taper, binning, spectral
model, or other rule; rank and positive-definiteness policy. \\
Output label & Second moment, randomization covariance, physical covariance,
marginal inverse variance, precision, effective consumer weight, formal
standardized signal, or empirical S/N. \\
\bottomrule
\end{longtable}

The word ``empirical'' identifies how a quantity was estimated; it does not
identify the target law or prove calibration.  The word ``weight'' identifies
neither marginal variance nor precision.  Those meanings must be explicit.

\section{Signal-to-noise and significance}

\subsection{Fixed scalar signal estimator}

Let $D$ be a science data vector on a valid domain, and let a fixed linear
signal estimator be
\begin{equation}
  \widehat a=\ell^{\mathsf T}D.
  \label{eq:linear-amplitude}
\end{equation}
For target covariance $T_0$ under a declared null hypothesis $H_0$, its noise
variance is
\begin{equation}
  \sigma_a^2=\ell^{\mathsf T}T_0\ell.
  \label{eq:linear-amplitude-variance}
\end{equation}
The standardized ratio
\begin{equation}
  Z_{\rm plug}=\frac{\widehat a-a_0}{\widehat\sigma_a}
  \label{eq:plugin-standard-score}
\end{equation}
is an S/N-like statistic.  It is statistical significance only after the
null law of the complete statistic is established.  If $T_0$ is known and
$\widehat a$ is Gaussian with null mean $a_0$, then $Z$ is standard normal.
If $T_0$ is estimated, the numerator shares data with the denominator, the
position or template was selected, or the conditional ensemble is non-Gaussian,
the standard-normal conclusion does not follow.  A Student-$t$, Hotelling,
random-field, bootstrap, randomization, or empirical-null calibration requires
its own exact assumptions; none is selected here.

For a fixed template $t$ and known positive-definite $T_0$, the generalized
least-squares amplitude is
\begin{equation}
  \widehat a_{\rm GLS}
  =\frac{t^{\mathsf T}T_0^{-1}D}
         {t^{\mathsf T}T_0^{-1}t},
  \qquad
  \Var_0(\widehat a_{\rm GLS})
  =\frac{1}{t^{\mathsf T}T_0^{-1}t}.
  \label{eq:gls-amplitude}
\end{equation}
This identity requires the template, response, support, and covariance to
refer to the same realized operator chain.  Plugging a noisy, singular, or
regularized covariance estimate into Equation~\eqref{eq:gls-amplitude}
changes the estimator and its null law.  The same ensemble cannot validate
that plug-in law merely because it supplied the covariance matrix.

The repository convention permits the truthful label
\texttt{formal\_standardized\_signal} for a signal multiplied by the square
root of a formal mapmaker weight when empirical noise products are unavailable.
That label is not statistical significance.  A product called
\texttt{sig2noise} requires an empirically calibrated uncertainty for the
exact numerator and response, plus the target and validity facts above.

\subsection{Thresholds, searches, and multiplicity}

A threshold belongs to a statistic and a decision loss, not to a map name.
For a fixed location, a one-sided threshold $q_\alpha$ may be defined by
\begin{equation}
  \Pr_{H_0}(Z\ge q_\alpha\mid\ThetaV,\mathcal V)\le\alpha.
  \label{eq:pointwise-threshold}
\end{equation}
A source finder usually searches positions, scales, templates, or iterations.
If
\begin{equation}
  Z_{\max}=\max_{j\in\mathcal J(D)}Z_j,
  \label{eq:search-maximum}
\end{equation}
then a global false-alarm threshold requires the null law of $Z_{\max}$,
including spatial covariance, search mask, boundaries, template bank,
selection rule, and any data-dependent set $\mathcal J(D)$.  Pointwise
Gaussian tail labels do not supply family-wise error, false-discovery rate,
completeness, or purity.  Threshold metadata must therefore name the statistic,
tail direction, null, search domain, multiplicity policy, calibration sample,
and version.

Source finding also conditions later photometry on a selection event.  Flux
and S/N reported at a detected maximum are subject to maximization and
selection bias.  A valid catalog claim must calibrate detection and estimation
jointly or separate discovery from estimation with an approved split or
cross-fitting design.

\subsection{Adaptive masks and feedback}

Suppose a thresholded source set or mask updates the next state according to
\begin{equation}
  \ThetaV_{k+1}
  =\Psi\!\left(\ThetaV_k,
    \mathcal F(D,\widehat T_k,q_k)\right),
  \label{eq:feedback-map}
\end{equation}
where $\mathcal F$ is a finder or classifier.  The R3 conditional covariance
at fixed $\ThetaV_k$ does not give the law of the adaptive sequence.  Noise
estimation, S/N, thresholding, masking/modeling, and the next noise estimate
form a closed statistical loop.  Valid feedback requires one of the following
independently approved designs:
\begin{itemize}
\item replay the complete adaptive procedure under a null or injected-signal
law, including stopping and every data-dependent operator;
\item estimate the noise and make feedback decisions from independent splits;
\item use cross-fitting with preserved fold identity and a proved aggregation
rule; or
\item establish another joint selective-inference law for the exact update.
\end{itemize}
The contract must separately define convergence, maximum iterations, failure,
oscillation, source preservation, and false-selection behavior.  An S/N value
alone is not a feedback or fruit-loop convergence criterion, and this core
does not authorize any FRUIT behavior.

\section{Aperture and other projected uncertainties}

Let $m\in\mathbb R^P$ be a map with target covariance $T$, and let
$a\in\mathbb R^P$ encode a fixed aperture, background subtraction, pixel-area
factor, mask, and any deterministic response normalization.  The extracted
quantity is
\begin{equation}
  \widehat f=a^{\mathsf T}m,
  \qquad
  \Var(\widehat f)=a^{\mathsf T}Ta.
  \label{eq:aperture-variance}
\end{equation}
Equivalently,
\begin{equation}
  a^{\mathsf T}Ta
  =\sum_pa_p^2T_{pp}
   +2\sum_{p<q}a_pa_qT_{pq}.
  \label{eq:aperture-cross-terms}
\end{equation}
Pixelwise variances are therefore insufficient unless off-diagonal terms are
proved irrelevant for that exact aperture and background operator.  A global
scale match is likewise insufficient when the correlation shape is wrong.

If $a^{\mathsf T}t$ is the response to a source template $t$, a response-
corrected flux and variance are
\begin{equation}
  \widehat F=\frac{a^{\mathsf T}m}{a^{\mathsf T}t},
  \qquad
  \Var(\widehat F)
  =\frac{a^{\mathsf T}Ta}{(a^{\mathsf T}t)^2},
  \label{eq:response-corrected-aperture}
\end{equation}
provided $a^{\mathsf T}t$ is fixed, finite, nonzero, and uses the same filter
and map response.  A fitted or data-selected aperture, centroid, radius,
background annulus, or template is random; its uncertainty requires the joint
procedure rather than Equation~\eqref{eq:aperture-variance} with the selected
operator treated as fixed.

An ensemble-based aperture estimator applies the exact same fixed $a$ to
every compatible realization and then estimates the scalar variance with the
finite-design rules of this document.  It is admissible only if R3 covariance
adequacy holds on the $a$ projection, source imprint and calibration are
controlled, the realized support is common, and the completed assignment
design gives adequate quadratic-statistic information.  No universal pixel
count, aperture size, or realization count follows.

\section{Use-specific adequacy hierarchy}

For each claimed use $u$, define a complete object
\begin{equation}
  \mathcal A_u=
  (\ThetaV,L_u,\mathcal V_u,\mathcal P_u^\star,
   \widehat T_u,\mathcal C_u,\mathcal D_R,\mathcal G_u),
  \label{eq:adequacy-object}
\end{equation}
where $\mathcal P_u^\star$ is the target law, $\mathcal D_R$ is the finite
joint assignment design, and $\mathcal G_u$ is the downstream decision rule.
The following levels are distinct:
\begin{enumerate}
\item \textbf{Identity adequate:} exact parent, operator, response, unit,
shape, indexing, support, validity, realization identity, and provenance match.
\item \textbf{Mean adequate:} the projected ensemble mean equals the target
mean to an owner-approved criterion, including any source-imprint condition.
\item \textbf{Scale adequate:} one predeclared scalar functional of projected
covariance matches the target.  This can support only that scale use.
\item \textbf{Marginal-variance adequate:} the required covariance diagonal
matches.  This does not establish off-diagonal covariance or precision.
\item \textbf{Covariance adequate:} $L_uK^\epsilon L_u^{\mathsf T}=T_u^\star$
on the complete use domain, or meets a preregistered approximation criterion.
\item \textbf{Precision adequate:} covariance estimation, rank, regularization,
and inversion error are acceptable for the exact inverse problem.
\item \textbf{Distribution/tail adequate:} the complete projected null law,
including estimator and selection, is adequate for the claimed tail or
interval probability.
\item \textbf{Adaptive-procedure adequate:} the full update, stopping, and
selection sequence has an accepted joint law and operating characteristics.
\end{enumerate}
Scale adequacy does not imply marginal adequacy; marginal adequacy does not
imply covariance or precision adequacy; covariance adequacy does not imply a
tail law; and a fixed-statistic tail law does not imply adaptive-procedure
adequacy.  Distributional adequacy implies lower moment adequacy only when the
required moments exist and the same conditioning domain is used.

The minimum level differs by use:
\begin{longtable}{>{\raggedright\arraybackslash}p{0.28\textwidth}
                  >{\raggedright\arraybackslash}p{0.65\textwidth}}
\toprule
Claimed use & Minimum independent requirement \\
\midrule
\endhead
Pixel variance diagnostic & Identity, declared centering/normalization, and
marginal randomization-variance estimation; no physical-noise claim implied. \\
Calibrated inverse-variance map & Marginal target adequacy, validated
calibration, inversion-bias treatment, and explicit nonprecision distinction. \\
Covariance-aware filter or GLS & Covariance and precision adequacy on the
filter/template subspace, including regularization and response. \\
Pixel S/N & Calibrated numerator variance for the exact pixel statistic;
truthful S/N label without an automatic significance claim. \\
Significance or source threshold & Distribution/tail adequacy plus null,
selection, search, and multiplicity calibration. \\
Aperture uncertainty & Covariance adequacy on the complete aperture,
background, and response projection; joint treatment if the aperture is fit. \\
Feedback or iterative masking & Adaptive-procedure adequacy or an approved
independence/cross-fitting design. \\
Production default & Every enabled consumer's level above, plus bounded Monte
Carlo uncertainty, resource/failure policy, and owner acceptance. \\
\bottomrule
\end{longtable}

\section{Production realization-count and default decisions}

No count is scientifically meaningful without a named target, assignment
design, estimator, consumer, and error criterion.  A production-count decision
must proceed in this order:
\begin{enumerate}
\item enumerate every enabled consumer and the adequacy level it claims;
\item fix the conditional ensemble, joint assignment design, target law,
centering, covariance estimator, calibration, and regularization;
\item select predeclared error metrics for the actual consumer functionals,
including bias, Monte Carlo dispersion, rank/conditioning, tail resolution,
and failure probability;
\item set owner-approved tolerances and confidence or coverage levels without
using the candidate count to choose a passing criterion;
\item derive or measure the use-specific information under the joint design,
including unique assignments and dependence of both linear and quadratic
summaries;
\item include requested, effective, completed, missing, duplicate, complement,
and use-specific effective counts in the decision;
\item specify resource, persistence, partial-failure, retry, and reproducibility
policy; and
\item validate the chosen fixed or adaptive stopping rule on representative
declared domains before making it a default.
\end{enumerate}

Different consumers can require different counts.  A single production stack
may support a pixel-scale diagnostic but not a covariance inverse, aperture,
or rare-tail threshold.  If one count is shared, its accepted scope is the
intersection of the uses it actually satisfies; it is not promoted by the
least demanding consumer.  A requested count is intent, an effective count is
the admitted plan, a completed count is realized state, and a unique or
effective count is a use-specific information descriptor.

Some estimator-specific algebraic boundaries are exact but are not universal
production minima:
\begin{itemize}
\item $R=1$ has zero empirically centered scatter and no centered sample-
variance degree of freedom, although a known-mean second moment can be formed;
\item the iid $R-1$ covariance requires $R\ge2$;
\item a generic centered $P$-dimensional covariance requires $R\ge P+1$ even
to have possible full rank;
\item the iid Gaussian mean-unbiased full inverse in
Equation~\eqref{eq:inverse-covariance-bias} requires $R>P+2$; and
\item shrinkage or a parametric covariance model can operate below those rank
bounds only by making additional assumptions and a different claim.
\end{itemize}

For iid Gaussian scalar variance, Equation~\eqref{eq:variance-relative-error}
implies the approximate planning relation
\begin{equation}
  R\gtrsim1+\frac{2}{\delta_v^2}
  \label{eq:variance-count-planning}
\end{equation}
for relative standard deviation at most $\delta_v$.  This is an analytic
benchmark only.  Correlated assignments, non-Gaussian fourth moments,
calibration uncertainty, source imprint, and spatial pooling can make the
actual requirement larger or qualitatively different.

Tail calibration has a separate information scale.  With $R$ independent
draws, an empirical exceedance frequency changes in increments of $1/R$.
If zero exceedances are observed for a fixed event with probability $p$, the
exact iid probability of that result is $(1-p)^R$; a one-sided level
$1-\alpha$ upper bound is
\begin{equation}
  p_{\rm upper}=1-\alpha^{1/R}.
  \label{eq:zero-exceedance-bound}
\end{equation}
Dependence reduces or otherwise changes this resolution.  A finite stack that
is useful for variance-scale validation can therefore be incapable of
calibrating a rare source-finding tail.

A count of 64 illustrates why a count cannot authorize itself.  Under the
ideal iid Gaussian scalar benchmark, its variance estimate has relative
standard deviation $\sqrt{2/63}\simeq0.178$ before calibration uncertainty.
Its raw empirical exceedance frequency moves in steps of $1/64$.  These facts
neither reject nor approve 64 for any use; they only show that 64 is optional
capacity whose adequacy must be assessed against a declared purpose.  It is
never presumed to be a production default, minimum, adequacy threshold, cap,
or universal requirement.

An adaptive count may stop when a preregistered precision diagnostic is met,
but the stopping statistic, minimum/maximum resource bounds, dependence-aware
confidence rule, failure route, and inference conditional on stopping must be
declared.  Stopping after inspecting a favorable weight, S/N, or tail result
without such a rule creates selection bias and cannot establish a default.

\section{Analytic limits}

Any later source trace or proportional fixture plan must recover the following
limits before an empirical or production interpretation is considered:
\begin{enumerate}
\item \textbf{Zero input.}  If every R3 realization is the zero vector, all
means, moments, scatters, projected variances, and aperture variances vanish;
no inverse weight is valid because the zero variance requires a declared
degenerate model rather than division by zero.

\item \textbf{One realization.}  $R=1$ gives $S_R=0$ identically.  A formula
that returns a positive centered empirical variance from one completed map is
using an unstated prior, pooling rule, known-mean moment, or external model.

\item \textbf{Independent common law.}  For iid realizations with estimated
mean, $S_R/(R-1)$ is unbiased for $K$ and $S_R/R$ is biased low by
$(R-1)/R$.  This is the only ordinary Bessel case.

\item \textbf{Known centered law.}  If $\mu=0$ is exact, the uncentered
$R^{-1}\sum_rY^{(r)}Y^{(r)\mathsf T}$ is unbiased for $K$ regardless of
cross-position dependence in its marginal expectation, but dependence still
controls its Monte Carlo error and deterministic source imprint can make the
target inappropriate for physical noise.

\item \textbf{Exact duplicates.}  If $Y^{(r)}=Y^{(1)}$ for every $r$, then
$S_R=0$ and the mean-specific information is at most one realization.  A
large completed count does not restore information.

\item \textbf{One complement pair.}  For $R=2$ with
$Y^{(2)}=-Y^{(1)}= -Y$, $\bar Y=0$ and
\begin{equation}
  S_2=2YY^{\mathsf T}.
  \label{eq:complement-pair-scatter}
\end{equation}
If the pair exactly enumerates a symmetric two-point marginal law, the
uncentered average $YY^{\mathsf T}$ is its covariance, whereas the iid
$R-1$ formula gives $2YY^{\mathsf T}$.  Complement pairing can make odd
moments exact while contributing only one distinct even-moment pattern.

\item \textbf{Exchangeable correlation.}  Common cross-covariance $C$ gives
Equation~\eqref{eq:exchangeable-scatter}; positive $C$ suppresses empirical
scatter and negative $C$ inflates it.  Unless $C$ is known and proportional to
$K$, a scalar effective-count correction is false.

\item \textbf{Rank.}  Centering $R$ maps supplies at most $R-1$ covariance
directions.  Any claimed full precision beyond that span must expose its
model or regularization.

\item \textbf{Linear projection.}  For fixed $L$,
\begin{equation}
  \Cov(LY)=LKL^{\mathsf T}.
  \label{eq:linear-projection-covariance}
\end{equation}
Two covariances with identical diagonals can give different aperture,
matched-filter, and smoothed-map variances.  A diagonal match cannot validate
those consumers.

\item \textbf{Scalar rescaling.}  If $T^\star$ is not proportional to
$\widetilde K$, no scalar $\alpha$ makes
$\alpha\widetilde K=T^\star$.  Global scale calibration cannot repair a
correlation-shape or morphology error.

\item \textbf{Source imprint.}  R3's deterministic-source term can elevate a
conditional second moment or covariance around source morphology even under a
centered sign law.  Inverting that elevation can create an apparent source-
shaped anti-weight.  Absence of the symptom does not prove source-free
physical covariance because filtering, normalization, support, finite design,
or calibration can hide it.

\item \textbf{Aperture correlation.}  If all aperture pixels have common
variance $v$ and common pairwise covariance $c$, an unweighted $n$-pixel sum
has variance
\begin{equation}
  nv+n(n-1)c,
  \label{eq:equicorrelated-aperture}
\end{equation}
not $nv$.  Even small $c$ can dominate a large aperture.

\item \textbf{Known Gaussian covariance.}  A fixed linear statistic divided
by its known null standard deviation is standard normal only under the stated
Gaussian null.  Estimating covariance, searching peaks, choosing apertures,
or feeding the result back changes that law.

\item \textbf{Complete assignment enumeration.}  If a deterministic design
exactly reproduces the target first and second sign moments, its conditional
mean and second moment can be exact without iid sampling.  Its uncertainty is
then design-specific; applying an iid finite-$R$ correction reintroduces error.
\end{enumerate}

\section{Falsification conditions and proportional future checks}

The following observations are sufficient to falsify the corresponding claim;
their absence is not proof of validity:
\begin{longtable}{>{\raggedright\arraybackslash}p{0.43\textwidth}
                  >{\raggedright\arraybackslash}p{0.50\textwidth}}
\toprule
Observation & Claim falsified \\
\midrule
\endhead
Requested $R$ differs from completed product identities without an explicit
partial-stack state & Any estimator normalized by requested count. \\
Duplicate or complement census conflicts with the declared joint design &
The reported independence, effective count, or finite correction. \\
Empirical scatter changes under permutation only because weights or support
are recomputed & A fixed-operator covariance identity. \\
Known-mean and empirically centered products share one unlabeled name & A
unique variance meaning. \\
Covariance estimate is not symmetric positive semidefinite on its claimed
complete-case domain & An unregularized covariance claim. \\
Claimed precision has rank exceeding the empirical span without a declared
model & A direct empirical-precision claim. \\
Calibrated scale agrees but aperture or filtered projections disagree &
Covariance adequacy beyond scalar scale. \\
Source-shaped second-moment or anti-weight morphology tracks injected or
known deterministic structure & A source-free-noise interpretation. \\
Nominal pointwise false-alarm rate changes with search area or template bank &
A global significance threshold. \\
Feedback result changes when noise estimation and selection are split & A
fixed-state or independence-based adaptive claim. \\
Empirical tail claim is finer than completed unique assignment resolution &
The claimed tail calibration. \\
\bottomrule
\end{longtable}

Proportional future checks, if separately authorized, need only tiny symbolic
or enumerated cases:
\begin{enumerate}
\item two scalar iid draws to distinguish divisors $R$ and $R-1$;
\item duplicate and complement pairs to verify
Equations~\eqref{eq:exchangeable-scatter} and
\eqref{eq:complement-pair-scatter};
\item a two-coordinate covariance with equal diagonals and opposite
off-diagonals to expose aperture and filter dependence;
\item a rank-deficient $P>R-1$ stack to expose undeclared precision;
\item a fixed compact deterministic component plus random residual to expose
source imprint before any inverse operation;
\item a fixed-location null versus a two-location maximum to expose the
look-elsewhere effect; and
\item a fixed aperture versus a data-selected aperture to expose selection
conditioning.
\end{enumerate}
These are specifications only.  No helper, schema, verifier, test, reduction,
or evidence request is created or executed by this freeze, and no numerical
tolerance is selected.

\section{Estimator and consumer provenance contract}

The requested, effective, observation-resolved, and realized lifecycle remains
one-way.  At minimum, interpretation requires the following logical facts;
this list does not mandate a new persistence schema:
\begin{longtable}{>{\raggedright\arraybackslash}p{0.18\textwidth}
                  >{\raggedright\arraybackslash}p{0.75\textwidth}}
\toprule
Layer & Required SCI-NOI-002 facts \\
\midrule
\endhead
Requested & Requested realization count and product/consumer scope; requested
variance, covariance, calibration, weight, S/N, threshold, aperture, or
feedback policy; disabled expert values retained as request, not execution. \\
Effective & Admitted target law and ensemble identity; estimator/version;
centering; finite normalization; weighting/missingness; covariance domain;
calibration and regularization policy; explicit unsupported combinations. \\
Observation-resolved & Exact observation/coadd/filter parent; map/group/Stokes
identity; signal unit, response, WCS/frame, indexing, shape, common support and
validity; fixed consumer projection, template, aperture/background, search
domain, and applicable calibration identity. \\
Realized design & Completed realization IDs and count; missing/failure state;
duplicates, complements, unique raw/even patterns; joint dependence law or
classified gap; per-pixel/pair counts if support varies; exact normalization
and any use-specific effective-information diagnostic. \\
Realized estimator & Reference/empirical mean; second moment versus covariance;
raw and calibrated covariance identities; calibration fit and uncertainty;
rank, eigen/support diagnostics, regularization, inverse correction; units and
validity. \\
Realized consumer & Numerator estimator and response; marginal/effective/full
weight meaning; S/N versus significance label; null and tail law; threshold,
search and multiplicity; aperture; feedback/stopping state; product inventory
and raw-parent identity. \\
\bottomrule
\end{longtable}

Routine products need not persist dense covariance matrices when a compact
model, operator, or reconstruction contract is sufficient for the claimed
use.  A digest verifies identity but does not reconstruct missing values.
Absence of exact reconstruction is a classified limitation, not permission to
infer a covariance or weight from a filename or unit.

\section{Decisions deliberately left open}

This core does not choose:
\begin{itemize}
\item whether any R3 conditional randomization product is adequate for a
repeated-physical-noise target;
\item known-zero versus empirical centering, divisor $R$ versus $R-1$, or a
joint-design correction;
\item a covariance representation, pooling domain, shrinkage, regularization,
or inverse correction;
\item calibration data, calibration operator, source-imprint treatment, or
acceptance tolerance;
\item marginal inverse variance, full precision, or consumer-effective weight
as the appropriate output;
\item any S/N numerator, null distribution, significance label, threshold,
search/multiplicity rule, source finder, or catalog policy;
\item any feedback, masking, cross-fitting, stopping, or convergence rule;
\item an aperture, background, source response, or adaptive photometry rule;
or
\item any production realization count or default.
\end{itemize}
Each decision requires a named owner, target, evidence route, failure policy,
and validation criterion.  No implementation or production disposition follows
from the equations alone.

\section{Phase boundary}

This file is the sole Phase 1 artifact.  The first return point is coordinator
verification and acceptance of its exact committed bytes.  Only a separate
Phase 2 authorization after that acceptance may open the sealed
\texttt{SCI-NOI-002-XAUD-001} and inspect the exact application source at
\texttt{d5015fe716971bf8ea617e8a187311bf5af05185}.  The first permitted source
boundary is the narrow realization-stack producer/persistence interface and
the first consumer of its variance/covariance/weight products, sufficient to
bind:
\begin{itemize}
\item completed count, centering, finite normalization, and cross-realization
design;
\item second-moment, covariance, calibration, weight, unit, support, validity,
and provenance labels;
\item the exact S/N, threshold, source-finding, feedback, and aperture
operators that are actually reachable; and
\item the production-count request/effective/realized path without assuming a
policy.
\end{itemize}
That trace must identify consumer symptoms without treating a consumer as
validation of its own estimator.  It remains documentation/audit work only.
No test, reduction, Unity action, evidence request, repair, integration, push,
or production authorization is implied.

\end{document}
